\newcommand{\ReplAbstractSentence}{All ten accuracy values of both released models were reproduced within 0.04 percentage points (vehicle mAP@0.5 78.00\% vs.\ 78.0\%); the reported parameter count of A-YOLOM(n) was not (3.64 vs.\ 4.43\,M), and Apple's MPS backend silently produced wrong metrics.}
\newcommand{\ReplDataNote}{The repository README swaps the download links of the drivable-area and lane annotations (the archive linked as \texttt{seg-drivable-10} contains the lane labels, class~11, and vice versa); we identified the archives by their class identifiers. The validation split contains 108{,}399 vehicle instances (car, bus, truck and train merged into one class, as in the original evaluation).}
\newcommand{\ReplEnvironment}{Apple M3~Pro (11-core CPU with 5 performance and 6 efficiency cores, 14-core GPU), 18\,GB unified memory, macOS~27.2.1; Python~3.10.22, PyTorch~2.2.2, torchvision~0.17.2, NumPy~1.26.4, OpenCV~4.9.0.80 and the ultralytics~8.0.105 fork contained in the repository (full list in \path{replication/requirements-lock.txt}). The official validator uses no dataloader worker processes on the CPU. One full evaluation took 51.8\,min for A-YOLOM(n) and 105.9\,min for A-YOLOM(s).}
\newcommand{\ReplDiscussion}{\textbf{Accuracy.} All ten accuracy values in Table~\ref{tab:repl} agree with the published ones within $\pm$0.04\,pp. The paper reports one decimal, so every difference is within the rounding interval ($\pm$0.05\,pp): the published evaluation of the released checkpoints is reproduced exactly, even with a newer PyTorch version (2.2.2 instead of 1.13.1) and on a CPU instead of an NVIDIA GPU. For completeness, the vehicle mAP@0.5:0.95 is 48.81\% for A-YOLOM(n) and 51.56\% for A-YOLOM(s).

\textbf{Aggregation of segmentation metrics.} The official code computes IoU and accuracy for every image and averages them. Accumulating one confusion matrix over the whole validation set gives values that are 2.1--2.8\,pp higher (Table~\ref{tab:agg}), because images with few lane or drivable-area pixels, where a few wrong pixels cost much, count as much as images with many such pixels. Neither choice is wrong, but the difference is larger than most improvements claimed in Table~\ref{tab:comparison}, and the aggregation is a detail of the evaluation code that readers cannot see unless the code is released. This gives a measured example of the comparability problem identified in RQ2.

\textbf{Parameters.} For A-YOLOM(s), the parameter count of the released checkpoint (13.61\,M before Conv+BN fusion) matches the paper. For A-YOLOM(n), both the checkpoint and the model built from the repository's configuration file have 3.64\,M parameters, not 4.43\,M as printed in the paper and README. The accuracy of this checkpoint matches the paper, so we consider the published parameter count a reporting error; its cause cannot be determined from the released material. After fusion, as used at inference, the models have 3.63\,M and 13.59\,M parameters.

\textbf{Speed.} On the M3~Pro GPU, A-YOLOM(n) reaches 37.3\,FPS at batch size~1 and 101.2\,FPS at batch size~32; A-YOLOM(s) reaches 29.5 and 53.5\,FPS. On the CPU (batch size~1) the values are 11.4 and 6.1\,FPS. Absolute values cannot be compared with the GTX~1080~Ti numbers of the paper, but the relative speed of the two models can: the paper reports a ratio n/s of 1.01 at batch size~1 and 1.79 at batch size~32; we measure 1.26 and 1.89. Both measurements time only the network forward pass, without pre-processing, NMS or mask post-processing, so they overestimate the end-to-end frame rate.

\textbf{Wrong metrics on the MPS backend.} Our first evaluation of A-YOLOM(n), with identical code, weights and data on the Apple GPU (MPS backend), gave a vehicle mAP@0.5 of 44.51\% instead of 78.00\% (Table~\ref{tab:mps}). A layer-wise comparison of the forward pass on one validation image showed no relevant difference between CPU and MPS (maximum relative difference $<3\cdot10^{-6}$). A minimal test (\path{replication/mps_inplace_view_repro.py}) shows that in PyTorch~2.2.2 an in-place \texttt{clamp\_} applied to a view of a tensor has no effect on MPS. The official post-processing and metric code use exactly this pattern: \texttt{clip\_boxes} clips predicted boxes to the image with \texttt{boxes[..., 0].clamp\_(0, w)}, and \texttt{update\_metrics\_seg} binarises masks with \texttt{mask\_list[i].clamp\_(max=1)}. We therefore attribute the discrepancy to this backend defect. No error or warning is raised, so without the published numbers as a reference this run could easily have been reported as a failed replication. Replication reports should therefore state the compute backend and, where possible, verify metrics on a second device.}
\newcommand{\ReplConclusion}{The replication confirms the published accuracy of A-YOLOM within 0.04\,pp for all ten metrics of both released models, but not the published parameter count of A-YOLOM(n), and it shows that the reported numbers also depend on how segmentation metrics are aggregated and on a correctly working compute backend}
